\documentclass[11pt]{article}
\usepackage[T1]{fontenc}
\usepackage{lmodern,amsmath,amssymb,amsthm,mathtools}
\usepackage[hmargin=0.90in,vmargin=0.80in]{geometry}
\usepackage[hidelinks]{hyperref}
\newtheorem{theorem}{Theorem}
\newtheorem{lemma}[theorem]{Lemma}
\newtheorem{corollary}[theorem]{Corollary}
\newcommand{\N}{N_\pi^m}

\title{Power sequences and delayed disagreement\\in consecutive permutation packings}
\author{}
\date{27 November 2026}
\hypersetup{pdftitle={Power sequences and delayed disagreement in consecutive permutation packings}}
\begin{document}
\maketitle
\begin{abstract}
We characterize every pattern and fixed overlap whose packing counts are powers of a single integer.
Two counts suffice, regardless of overlap size. For $m\ge2$, copies two and $m+1$ suffice.
In the power classes, the packing poset splits into singletons and pairs of disjoint chains.
For base two, each overlap contributes exactly one incomparable pair.
We also construct pairs of patterns whose counts first differ arbitrarily late.
Consequently, no absolute copy cutoff determines Wilf equivalence.
Lean 4 verifies the counterexample to an absolute cutoff.
These results address Questions 8 and 9 of Blitvi\'c, Kammoun, and Steingr\'imsson.
\end{abstract}

\section{The classification}
Fix a permutation $\pi$ of $[K]=\{1,\ldots,K\}$ and an overlap $1\le m<K$.
Write $d=K-m$. For $j\ge1$, take $j$ windows of length $K$ with starts $0,d,\ldots,(j-1)d$.
Let $\N(j)$ count permutations of $[K+(j-1)d]$ that induce $\pi$ in every prescribed window.
Additional occurrences are allowed. Two pattern--overlap pairs are Wilf equivalent if their counts agree for every $j$ \cite{bks}.
We ask when $\N(j)=k^{j-1}$ for a positive integer $k$.

Admissibility means that the first and last $m$ entries have the same relative order.
Failure of admissibility gives $\N(2)=0$. Assume admissibility throughout.
Let $p_1<\cdots<p_m$ and $q_1<\cdots<q_m$ list the prefix and suffix ranks, respectively.
The overlap identifies rank $q_i$ on the left with rank $p_i$ on the right.
Write $P=\{p_i\}$ and $Q=\{q_i\}$. Set $p_0=q_0=0$ and $p_{m+1}=q_{m+1}=K+1$.
Define
\[
 a_i=p_{i+1}-p_i-1,\qquad c_i=q_{i+1}-q_i-1,\qquad
 b=\prod_{i=0}^m\binom{a_i+c_i}{a_i}.
\]
For $z\in\{0,\ldots,K+1\}$ define two nondecreasing maps
\begin{align*}
 D(z)&=\max\bigl(\{0\}\cup\{p_i:q_i\le z,\ 1\le i\le m\}\bigr),\\
 U(z)&=\min\bigl(\{K+1\}\cup\{p_i:q_i\ge z,\ 1\le i\le m\}\bigr).
\end{align*}

\begin{theorem}\label{class}
For every admissible overlap, the following conditions are equivalent.
\begin{enumerate}
\item $\N(j)=k^{j-1}$ for every $j\ge1$.
\item $\N(2)=k$ and $\N(m+2)=k^{m+1}$.
\item $b=k$, and for every $x\in[K]\setminus Q$, $y\in[K]\setminus P$, and $2\le t\le m+1$,
\[
 y\le D^t(x)\quad\text{or}\quad y\ge U^t(x).
\]
\end{enumerate}
\end{theorem}

\begin{proof}
Let the order of each window follow the ranks of $\pi$.
Their transitive closure is the packing poset. Sorting positions by value identifies packings with linear extensions \cite[Theorem~5.3]{bks}.
Write $W_i$ for window $i$ and $B_i=W_i\cup W_{i+1}$ for an adjacent pair.
Inside $B_i$, the shared chain separates $m+1$ gaps.
Each gap consists of two disjoint chains of lengths $a_h,c_h$.
Their interleavings give the binomial factors in $b$. Hence $\N(2)=b$.

Choose one total order on each $B_i$ that respects both windows.
These choices extend to a global order.
Indeed, $B_i$ meets the preceding union exactly in $W_i$. The two orders agree on that chain.
Between consecutive shared elements, merge the two private lists while preserving each list's order.
Merge the lists at the two ends in the same way.
This proves surjectivity of restriction onto the $b^{j-1}$ tuples. Hence $\N(j)\ge b^{j-1}$.
The same construction extends every order on a consecutive block by adjoining windows on either side.
Thus later attachments preserve the induced order on every such block.
Borga and Penagui\~ao use related gluing for compatible consecutive patterns \cite{bp}.

We claim equality holds exactly when every incomparable pair lies in some $B_i$, for $j\ge2$.
Under this condition the tuple determines every pairwise comparison.
Conversely, take incomparable positions $u,v$ outside any common $B_i$.
The union of their strict predecessor sets forms an order ideal containing neither position.
List that ideal first, then $u,v$, then the remaining positions in any compatible order.
Swapping $u,v$ gives a different order with the same tuple. This proves the claim.

The windows containing a position form an interval of indices.
Two positions share no $B_i$ exactly when their index intervals have distance at least two.
For such positions, let $r$ end the earlier interval and let $r+t$ start the later interval.
Their endpoint ranks satisfy $x\notin Q$ and $y\notin P$.
Conversely, each such pair of ranks occurs at these endpoints in a block of $t+1$ windows.
We may compute comparisons inside this block, by the extension argument.

At distance $t$, the ranks weakly below the initial position $x$ are exactly the ranks at most $D^t(x)$.
The ranks weakly above it are exactly the ranks at least $U^t(x)$.
At $t=0$, both boundaries equal $x$.
For the induction step, the incoming window meets the earlier union exactly in its intersection with the preceding window.
Indeed, an old position in the incoming window lies after its start and before the preceding window's end.
It therefore belongs to both windows, even if three or more windows contain it.

A comparison crossing between old and new positions must pass through this shared chain.
A path returning to the shared chain follows its fixed order and creates no further comparisons.
If the previous lower boundary is $z$, the admissible shared ranks are precisely $q_i\le z$.
Their new ranks are the corresponding $p_i$.
Monotonicity puts a shared position of new rank $p_s$ below $x$ exactly when $q_s\le z$.
Taking the largest available $p_i$, or zero, gives $D(z)$.
Reversing the inequalities gives $U(z)$. This completes the induction.

Consequently, a strict gap $D^t(x)<y<U^t(x)$ at $t\ge2$ witnesses failure of equality.
An orbit of a nondecreasing map on a finite chain moves in one direction.
After one step, each map takes at most $m+1$ values.
Each orbit must therefore stabilize by step $m+1$.
Every possible failure occurs within the tested distances.
A failing block extends to $m+2$ windows while preserving its incomparable pair.
Its endpoint positions still share no $B_i$ because $x\notin Q$ and $y\notin P$.
This proves the equivalences. One window always has one packing.
\end{proof}

\begin{corollary}[A shorter test]\label{short}
Set $J=\max(3,m+1)$. Then $\N(j)=k^{j-1}$ for every $j\ge1$ if and only if
$\N(2)=k$ and $\N(J)=k^{J-1}$.
\end{corollary}
\begin{proof}
Theorem~\ref{class} gives the result for $m=1$.
For $m\ge2$, suppose the first strict gap occurs at distance $m+1$.
At distance $m$, the same endpoint rank $y$ lay outside the gap.
Hence the lower boundary falls past $y$ or the upper boundary rises past $y$ at the last step.
Rank complementation exchanges these cases, so consider the falling lower boundary.
If two consecutive iterates agree, all later iterates agree.
Thus $D(x),\ldots,D^{m+1}(x)$ decrease strictly through all $m+1$ possible values.
Their last two values are $p_1$ and $0$. Hence $1\le y<p_1$, since $y\notin P$.
Also $D(p_1)=0$ gives $p_1<q_1$. Consequently $1\notin Q$, $D^2(1)=0$, and $U^2(1)=p_1$.
The same $y$ already lies in a strict gap at distance two, contradicting first failure.
Thus distances $2,\ldots,m$ suffice. The block-extension argument gives the improved copy cutoff.
\end{proof}
For $m=1$, three copies are necessary in general.
For instance, $132$ has three two-copy packings and fifteen three-copy packings.

\paragraph{Chain case.}
At $k=1$, Theorem~\ref{class} characterizes the pairs whose poset is a chain for every $j$.
A finite poset has exactly one extension precisely when it is a chain.
Theorem~5.23 of \cite{bks} gives a gap-vector criterion for $m\le K/2$.
Theorem~\ref{class} covers $m>K/2$ as well.
For large overlaps, $b=1$ alone is insufficient, as Section~\ref{delay} shows.

\section{Structure at equality}
\begin{theorem}\label{blocks}
Suppose $\N(j)=b^{j-1}$ for a fixed $j\ge2$.
The packing poset is an ordinal sum of singleton blocks and disjoint unions of two nonempty chains.
In an ordinal sum, every element of an earlier block precedes every element of a later block.
Each $B_i$ contributes one two-chain block of lengths $a_h,c_h$ for each gap with $a_hc_h>0$.
\end{theorem}

\begin{proof}
Join two positions by an edge when they are incomparable.
Every edge belongs to some $B_i$ by the equality criterion.
For $h>i$, $B_i\cap B_h$ lies in the chain $W_{i+1}$.
Thus each edge belongs to exactly one $B_i$. Label it by $i$.

Suppose $uv$ and $uw$ have different labels.
If $v<w$, choose local orders imposing $w<u<v$.
Independent gluing realizes both choices, contradicting $v<w$.
The same argument excludes $w<v$. Thus $v,w$ form another edge.
Its label differs from both previous labels. Otherwise one double window would contain all three endpoints and repeat another edge.
Orient the three edges cyclically in their three local orders. Gluing now gives a contradiction.
Therefore all edges incident at one vertex have the same label.
Following paths proves that each connected component with edges has one label.

Within $B_i$, the two chains of each nonempty mixed gap form a complete bipartite incomparability graph.
Distinct gaps have no edges between them. Extension preserves these comparisons and incomparabilities.
The label result therefore identifies all global components.
An outside vertex compares with every vertex in a component.
It lies on the same side of both ends of any edge, by transitivity.
Following paths in each of two components orders the components uniformly. Their order gives the asserted ordinal sum.
\end{proof}

For $b=2$, the product has exactly one nontrivial factor, $\binom21$.
Hence exactly $j-1$ blocks are isolated incomparable pairs. Every other block is a singleton.
Conversely, these blocks give $2^{j-1}$ extensions.
This proves the base-two observation in \cite[Section~7, item~9]{bks}.
A separator in the sense of item~5 is a singleton with at least one two-chain block on each side.
This describes separators within these power classes.

\paragraph{Interleaving ranks.}
Take $123584679$ with $m=4$. Here $P=\{1,2,3,5\}$ and $Q=\{4,6,7,9\}$.
The only mixed gap has lengths one and one, so $b=2$.
All tested $x$ satisfy $x\le8$, and all tested $y$ satisfy $y\ge4$.
The upper orbit from $8$ runs through $5,2,1,1,\ldots$.
Hence $U^t(x)\le2<y$ throughout the test.
By contrast, $1324$ with $m=2$ has $b=2$ and $\N(3)=5$.

\paragraph{Infinite classes.}
For any $k\ge1$, take $K\ge k+4$, prefix $(1,k+1)$, and suffix $(K-2,K)$.
Fill the middle with the unused ranks in increasing order. Take $m=2$.
Each shared interval interleaves one left element with $k-1$ ordered right elements.
In each interior window, the first shared interval ends before the last shared interval begins in value order.
All other comparisons are forced by the window chains.
Thus the $j-1$ insertions are independent and give exactly $k^{j-1}$ extensions.
Varying $K$ gives infinitely many representatives for every positive base.

\section{No absolute cutoff for Wilf equivalence}\label{delay}
Item~8 of \cite{bks} asks whether five copies distinguish all inequivalent pairs.
The following family rules out every constant bound, even for two patterns of equal length and equal overlap.

\begin{theorem}\label{delayed}
For every $L\ge5$, set $A=L-2$, $K=3L$, and $m=3L-3$.
Construct $\pi_L$ by replacing the following list with its ranks:
\[
 (n,\ n+A+\tfrac13,\ n+2A+\tfrac23),\qquad n=0,\ldots,L-1,
\]
with triples in increasing index order.
For every $j\ge1$,
\[
 N_{\pi_L}^m(j)=2^{\max(0,\,j-L+2)} .
\]
\end{theorem}

\begin{proof}
The stride is three, so a packing with $j$ windows uses triples $X_i,Y_i,Z_i$ for $0\le i\le M=L+j-2$.
Each prescribed window contains $L$ consecutive triples. Shifting all triple indices adds a constant to the displayed values.
Thus each such window has the same pattern.
Any indices at distance at most $L-1$ share a window.
Indeed, for $i\le h$ choose its start $\max(0,h-L+1)$.
Each of the three individual tracks increases.

For adjacent tracks we have
\[
 X_i<Y_h\ \Longleftrightarrow\ i\le h+A,\qquad
 Y_i<Z_h\ \Longleftrightarrow\ i\le h+A .
\]
For the first equivalence, if $i\le h$, use $X_i\le X_h<Y_h$.
If $h<i\le h+A$, compare directly in one window.
If $i\ge h+A+1$, compare $Y_h<X_{h+A+1}\le X_i$.
These direct comparisons span at most $A+1=L-1$ indices.
The same proof applies to $Y,Z$.

If $i\le h+2A$, then $t=\max(0,i-A)$ gives $X_i<Y_t<Z_h$.
If $i\ge h+2A+2$, then $t=h+A+1$ gives $Z_h<Y_t<X_i$.
Both intermediate indices lie between zero and $M$. Only the pairs $X_{h+2A+1},Z_h$ remain.

For $r\in\{1/3,2/3\}$ assign values $(n,n+A+r,n+2A+2r)$ to each triple.
Each assignment gives distinct values. Both induce the prescribed order in every window.
Indeed, only an outer-track comparison at index difference $2A+1$ can change, and $2A+1>L-1$.
These orders reverse every remaining pair, proving its incomparability.
The incomparable pairs have disjoint endpoints.
By transitivity, every third element lies below both endpoints or above both.
Hence each pair forms a two-element block, and its order can be chosen independently.
There are $\max(0,M-2A)$ such pairs. Substitution gives the formula.
\end{proof}

The increasing pattern of the same length has one packing for every $j$.
The first disagreement occurs at $j=L-1=m/3$. In particular, $L=7$ gives the pattern
\[
 (1,7,14,2,9,16,3,10,17,4,11,18,5,12,19,6,13,20,8,15,21).
\]
At overlap 18 its counts start $1,1,1,1,1,2,4,8$.
The family gives a linear lower bound along $m=3L-3$.
The shorter test gives the upper bound $\max(3,m+1)$ for comparison with a power sequence.
The best overlap-dependent cutoff remains undetermined in general.
The theorem also leaves open an upper bound of practical size for two arbitrary patterns of bounded length.

\paragraph{Delayed disagreement at each fixed base.}
For completeness, every fixed $k\ge2$ also admits arbitrarily delayed disagreement with $k^{j-1}$.
Use two groups of three tracks per time step.
For a three-track group, allow successive integral offsets $a,b$ with $1\le a,b\le L-2$ and $a+b+1>L-1$.
Its coordinates are $(n,n+a+r,n+a+b+2r)$.
The preceding proof applies with thresholds $a,b$ on the two adjacent track pairs.
For the forward outer comparison use $t=\max(0,i-a)$.
For the reverse comparison use $t=h+b+1$.
The only free pairs are $X_{h+a+b+1},Z_h$, giving $\max(0,M-a-b)$ isolated pairs.

Give the lower group offsets $(1,L-2)$ and reverse its value order.
Give the upper group offsets $(L-2,L-2)$ in the original order.
At each time, list the lower triple followed by the upper triple.
Place every lower value below every upper value in the pattern.
This comparison holds globally because each lower track decreases and each upper track increases.
Any lower position is at most its time-zero position.
That position precedes the upper time-zero position, which is at most the chosen upper position.
Thus the two groups form an ordinal sum.
The lower group contributes $j-1$ free pairs. The upper group contributes $\max(0,j-L+2)$.

In every lower triple, replace $X$ by $k-1$ consecutive positions with consecutive increasing ranks.
Replace that position consistently in every window. Each original forced comparison lifts to every member of the replacement chains.
Conversely, projecting an increasing path between distinct chains gives a path in the original order after removing repeated vertices.
Thus each lower free pair becomes a chain of length $k-1$ beside a singleton.
It contributes $k$ extensions. The full count is
\[
 k^{j-1}2^{\max(0,j-L+2)}.
\]
Replacing the upper group by a fully increasing list of triples gives a baseline pattern with counts $k^{j-1}$.
Both patterns have $K=(k+4)L$ and $m=(k+4)(L-1)$. They first differ at $j=L-1$.

\section*{Verification and provenance}
The accompanying Python program counts extensions from the original window inequalities using order ideals.
For lengths two through six it checks 5,154 counts across 1,462 admissible pairs.
At length seven it checks all 172 admissible pairs with overlaps four, five, or six.
For these pairs it checks 1,054 counts and 690 posets at equality.
It also checks the delayed family at eleven lengths from 15 through 45.
Nagy gives a Lean formalization for the double-factorial family in item~6 \cite{nagy}. Part~II develops other product formulas \cite{bks2}.
Lean checks the counterexample to every absolute cutoff using actual lists of positions ordered by rank.
Its theorem \texttt{Packing.Delay.no\_absolute\_wilf\_cutoff} compares the finite packing sets by bijection.
Lean also checks gluing, recovery from local orders, rank transfer, and the component argument.
The complete power-class classification and the exact counts after first disagreement have written proofs.
\paragraph{Declarations.}
Astra 6.2 produced this work. I selected the problems and released this work on my GitHub repository. I wrote none of the proofs. \href{https://github.com/palvimaki/permutation-power-classes/issues}{Contact me for corrections.}
Gemini Neon checked the proofs in the original power-class paper.

\small
\begin{thebibliography}{9}
\setlength{\itemsep}{1pt}
\raggedright
\bibitem{bks} N. Blitvi\'c, M. S. Kammoun, and E. Steingr\'imsson.
\emph{Permutations, Pattern Covers, Posets and Positivity: Part I: Patterns and Posets}.
\href{https://arxiv.org/abs/2610.29846}{arXiv:2610.29846}, 2026.
\bibitem{bp} J. Borga and R. Penagui\~ao.
\emph{The feasible region for consecutive patterns of permutations is a cycle polytope}.
\href{https://arxiv.org/abs/2003.12661}{arXiv:2003.12661}, 2020.
\bibitem{bks2} N. Blitvi\'c, M. S. Kammoun, and E. Steingr\'imsson.
\emph{Permutations, Pattern Covers, Posets and Positivity, Part II: Regular Pattern Packings and Moment Sequences}.
\href{https://arxiv.org/abs/2611.12829}{arXiv:2611.12829}, 2026.
\bibitem{nagy} Z. Nagy.
\emph{A uniform double-factorial construction for consecutive-pattern packings}.
\href{https://arxiv.org/abs/2611.06937}{arXiv:2611.06937}, 2026.
\end{thebibliography}
\end{document}
